\RequirePackage[thinlines]{easybmat} %-- muss aufgrund eines Bugs vor etex und tikz geladen werden
\documentclass[a4paper, twoside, headsepline, index=totoc,toc=listof, fontsize=10pt, cleardoublepage=empty, headinclude, DIV=13, BCOR=13mm, titlepage,]{scrartcl}
\usepackage{scrtime} % KOMA, Uhrzeit ermoeglicht
\usepackage{scrpage2} % wie fancyhdr, nur optimiert auf KOMA-Skript, leicht andere Syntax
\usepackage{etoolbox}
\usepackage{letltxmacro}

% Compiler abhaengige IFs
\usepackage{ifxetex,ifluatex}
\newif\ifxetexorluatex
\ifxetex
  \xetexorluatextrue
\else
  \ifluatex
    \xetexorluatextrue
  \else
    \xetexorluatexfalse
  \fi
\fi

%--Farbdefinitionen
%-- muss vor tikz geladen werden
\usepackage[usenames, table, x11names]{xcolor}
\definecolor{dark_gray}{gray}{0.45}
\definecolor{light_gray}{gray}{0.6}

%--Zum Zeichnen
%-- muss vor polyglossia bzw. babel geladen werden
\usepackage{tikz}
\usepackage{tikz-cd}
\usetikzlibrary{external}
\tikzset{>=latex}
\usetikzlibrary{shapes,arrows,intersections}
\usetikzlibrary{calc,3d}
\usetikzlibrary{decorations.pathreplacing,decorations.markings}
\usetikzlibrary{angles}

%-- Konfiguration von tikzexternalize
\tikzexternalize[prefix=tikz/, up to date check=diff]

\ifxetexorluatex
\tikzset{external/system call={xelatex \tikzexternalcheckshellescape %-- verwende LuaLaTeX, wegen dynamischer Speicherallokation
    -halt-on-error -interaction=batchmode -jobname "\image" "\texsource"}}
\else
\tikzset{external/system call={pdflatex \tikzexternalcheckshellescape 
    -halt-on-error -interaction=batchmode -jobname "\image" "\texsource"}}
\fi

%-- tikzexternalize fuer tikzcd deaktivieren, da inkompatibel
\AtBeginEnvironment{tikzcd}{\tikzexternaldisable}
\AtEndEnvironment{tikzcd}{\tikzexternalenable}

%-- um Inkompatibilitaeten von quotes und polyglossia bzw. babel zu vermeiden
\tikzset{
  every picture/.append style={
    execute at begin picture={\shorthandoff{"}},
    execute at end picture={\shorthandon{"}}
  }
}
\usetikzlibrary{quotes}
\usepackage{pgfplots}
\usepgfplotslibrary{colormaps}
\pgfplotsset{compat=1.10}



%-- Mathe
\usepackage{mathtools} % beinhaltet amsmath
\mathtoolsset{showonlyrefs, centercolon}
\newtagform{brackets}[\textbf]{[}{]}
\usetagform{brackets}
\usepackage{wasysym}
\usepackage{amssymb} %zusätzliche Symbole
\usepackage{latexsym} %zusätzliche Symbole
\usepackage{stmaryrd} %für Blitz
\usepackage{nicefrac} %schräge Brüche
\usepackage{cancel} %Befehle zum Durchstreichen
\usepackage{mathdots}
% \usepackage{bm}
\DeclareSymbolFont{bbold}{U}{bbold}{m}{n}
\DeclareSymbolFontAlphabet{\mathbbold}{bbold}
\newcommand{\mathds}[1]{\mathbb{#1}} %Um Kompatibilitaet mit frueheren Benutzung von dsfont herzustellen
\newcommand{\ind}{\mathbbold{1}} % charakteristische-Funktion-Eins
\def\mathul#1#2{\color{#1}\underline{{\color{black}#2}}\color{black}} %farbiges Untersteichen im Mathe-Modus

%-- Alles was mit Schrift und XeteX zu tun hat
\ifxetexorluatex
    % XeLaTeX or LuaTeX
	\ifxetex
		\usepackage{mathspec} %beinhaltet fontspec 
	\else
		\usepackage[no-math]{fontspec}
	\fi
    \usepackage{polyglossia} %babel-ersatz
    \setmainlanguage[spelling=new,babelshorthands=false]{german}
	\newcommand\glqq{"}
	\newcommand\grqq{"}
	\defaultfontfeatures{Mapping=tex-text, WordSpace={1.4}} %
    \setmainfont[Ligatures=Common, BoldFont={* Bold}, ItalicFont={* Light Italic},SmallCapsFont={Linux Libertine O}, SmallCapsFeatures={Letters=SmallCaps}]{Source Sans Pro}
	\setsansfont[Scale=MatchLowercase,Ligatures=Common, BoldFont={* Medium}]{Ubuntu}
	\ifxetex
		\setallmonofonts[Scale=MatchLowercase, ItalicFont={*}]{Consolas} 
	\else
		\setmonofont[Scale=MatchLowercase, ItalicFont={*}]{Consolas}
	\fi
	\usepackage{xltxtra}
	\usepackage{fontawesome}
	\usepackage{microtype}
\else
    % default: pdfLaTeX
    \usepackage[ngerman]{babel}
    \usepackage[T1]{fontenc}
    \usepackage[utf8]{inputenc}
	\usepackage{textcomp} %verhindert ein paar Fehler bei den Fonts
    \usepackage[babel=true, tracking=true]{microtype}
	\usepackage{lmodern}
	\renewcommand{\familydefault}{\sfdefault} 
\fi


%--Mixed
\usepackage[neverdecrease]{paralist}
\usepackage[german=quotes]{csquotes}
\usepackage{makeidx}
\usepackage{booktabs}
\usepackage{wrapfig}
\usepackage{float}
\usepackage[margin=10pt, font=small, labelfont={sf, bf}, format=plain, indention=1em]{caption}
\captionsetup[wrapfigure]{name=Abb. }
\usepackage{stackrel}
\usepackage{ifthen}
\usepackage{multicol}
\usepackage{qrcode}
\flushbottom

%--Hyphenation
\hyphenation{di-men-si-o-na-le}


%--Unterstreichung
\usepackage[normalem]{ulem}
\setlength{\ULdepth}{1.8pt}

%--Indexverarbeitung
\newcommand{\bet}[1]{\textbf{#1}}
\newcommand{\Index}[1]{\textbf{#1}\index{#1}}
\makeindex
\setindexpreamble{{\noindent \itshape Die \emph{Seitenzahlen} sind mit Hyperlinks zu den entsprechenden Seiten versehen, also anklickbar} \ifxetexorluatex \faHandUp\fi \par \bigskip}
\renewcommand{\indexpagestyle}{scrheadings}


%--Marginnote & todonotes
\usepackage{marginnote}
\renewcommand*{\marginfont}{\color{dark_gray} \itshape \footnotesize }
\usepackage[textsize=small]{todonotes}

%--Todonotes mit tikz/externalize kompatibel machen
\LetLtxMacro{\oldtodo}{\todo}
\renewcommand{\todo}[2][]{\tikzexternaldisable\oldtodo[#1]{#2}\tikzexternalenable}
\LetLtxMacro{\oldmissingfigure}{\missingfigure}
\renewcommand{\missingfigure}[2][]{\tikzexternaldisable\oldmissingfigure[{#1}]{#2}\tikzexternalenable}

%--Konfiguration von Hyperref pdfstartview=FitH, 
\usepackage[hidelinks, pdfpagelabels,  bookmarksopen=true, bookmarksnumbered=true, linkcolor=black, urlcolor=SkyBlue2, plainpages=false,pagebackref, citecolor=black, hypertexnames=true, pdfauthor={Jannes Bantje}, pdfborderstyle={/S/U}, linkbordercolor=SkyBlue2, colorlinks=false]{hyperref}

% \let\orighref\href
\ifxetexorluatex
\newcommand{\hrefsym}[2]{\href{#1}{\texttt{#2\,\faExternalLink}}}
\newcommand{\hrefsymmail}[2]{\href{#1}{\texttt{\faEnvelope\,#2}}}
\renewcommand{\url}[1]{\hrefsym{#1}{\nolinkurl{#1}}}
\fi
\providecommand{\hrefsym}{\href}
% \usepackage{cleveref}
% \crefformat{footnote}{#2\footnotemark[#1]#3}

%--Römische Zahlen
\newcommand{\RM}[1]{\MakeUppercase{\romannumeral #1{}}}

%--Punkte (nach hyperref)
\usepackage{ellipsis}

%%--Abkürzungen etc.
\newcommand{\light}{\text{\Large $\lightning$}}

%-- Definitionen von weiteren Mathe-Befehlen
\DeclareMathOperator{\re}{Re} %Realteil
\DeclareMathOperator{\im}{Im} %Imaginaerteil
\DeclareMathOperator{\id}{id} %identische Abbildung
\DeclareMathOperator{\Sp}{Sp} %Spur
\DeclareMathOperator{\sgn}{sgn} %Signum
\DeclareMathOperator{\alt}{Alt} %Alternierende n-linearFormen
\DeclareMathOperator{\End}{End} %Endomorphismen
\DeclareMathOperator{\Vol}{Vol} %Volumen
\DeclareMathOperator{\dom}{dom} %Domain
\DeclareMathOperator{\Hom}{Hom} %Homomorphismen
\DeclareMathOperator{\bild}{Bild} %Bild
\DeclareMathOperator{\Kern}{Kern }
\DeclareMathOperator{\rg}{rg} %Rang
\DeclareMathOperator{\Rg}{Rg} %Rang
\DeclareMathOperator{\diam}{diam} %Durchmesser
\DeclareMathOperator{\dist}{dist} %Distanz
\DeclareMathOperator{\grad}{grad} %Gradient
\DeclareMathOperator{\dive}{div} %Gradient
\DeclareMathOperator{\rot}{rot} %Rotation
\DeclareMathOperator{\hess}{Hess} %Hesse-Matrix
\DeclareMathOperator{\koker}{Koker} %Kokern
\DeclareMathOperator{\aut}{Aut} %Automorphismen
\DeclareMathOperator{\ord}{ord} %Ordnung
\DeclareMathOperator{\ggT}{ggT} %ggT
\DeclareMathOperator{\kgV}{kgV} %kgV
\DeclareMathOperator{\Gr}{Gr} %Gerade
\DeclareMathOperator{\Kr}{Kr} %Kreis
\DeclareMathOperator{\Char}{char} %Charakteristik
\DeclareMathOperator{\Aut}{Aut} %Automorphismen
% \DeclareMathOperator{\D}{D} %formale Ableitung
\newcommand{\D}{\ensuremath{\mathrm{D}\mkern-1.5mu}}
\newcommand{\mathd}{\ensuremath{\mathrm{d}\mkern-1.0mu}}
\newcommand{\Tmap}{\ensuremath{\mathrm{T}\mkern-1.0mu}}
\newcommand{\opL}{\ensuremath{\mathrm{L}\mkern-0.6mu}}
\DeclareMathOperator{\cov}{cov} %Kovarianz
\DeclareMathOperator{\Gal}{Gal} %Galois
\DeclareMathOperator{\supp}{supp} %Träger
\DeclareMathOperator{\Sym}{Sym} %Symmetrische Gruppe
\DeclareMathOperator{\Zyl}{Zyl} %Zylinder
\DeclareMathOperator{\Mod}{Mod} %Moduln
\DeclareMathOperator{\EPK}{EPK} %Einpunktkompaktifizierung
\DeclareMathOperator{\conj}{conj} %Konjugation
\DeclareMathOperator{\ann}{ann} %Annulator
\DeclareMathOperator{\tor}{tor} %Torsionsmodul
\DeclareMathOperator{\Int}{Int} %Interior/Inneres
\DeclareMathOperator{\Br}{Br} %Brauerergruppe
\DeclareMathOperator{\GL}{GL} %allgemeine lineare Gruppe
\DeclareMathOperator{\SL}{SL} %spezielle lineare Gruppe
\DeclareMathOperator{\SU}{SU} %spezielle unitäre Gruppe
\DeclareMathOperator{\Tr}{Tr} %Spur/Trace
\DeclareMathOperator{\conv}{conv} %Konvexe Teilmenge
% \DeclareMathOperator{\tr}{tr}
\newcommand{\tr}{\mathrm{tr}}
\newcommand{\w}{\mathrm{w}}
\newcommand{\weakT}[1]{\ensuremath{\mathcal{T}_{#1}^{\mkern+1.0mu\text{\raisebox{0.4ex}{$\mathrm{w}$}}}}}
\newcommand{\weakTstar}[1]{\ensuremath{\mathcal{T}_{#1}^{\mkern+1.0mu\text{\raisebox{0.4ex}{$\mathrm{w}$}}^*}}}
\newcommand{\rand}[1]{\ensuremath{\partial^{\scriptscriptstyle #1}}}


% -- Kategorien
\DeclareMathOperator{\Ob}{Ob} %Objekte
\DeclareMathOperator{\Mor}{Mor} %Morphismen

% \DeclareMathOperator{\TOP}{TOP} %Topologische Räume
\newcommand{\TOP}{\textsc{Top}}
% \DeclareMathOperator{\HTOP}{HTOP} %Topologische Räume + Homotopieäquivalenz
\newcommand{\HTOP}{\textsc{HTop}}
% \DeclareMathOperator{\VR}{VR} %Vektorräume
\newcommand{\VR}{\textsc{VR}}
% \DeclareMathOperator{\MOD}{MOD} %Moduln
\newcommand{\MOD}{\textsc{Mod}}
% \DeclareMathOperator{\MENGEN}{MENGEN} %Mengen
\newcommand{\MENGEN}{\textsc{Mengen}}
% \DeclareMathOperator{\MAN}{MAN} %Mannigfaltigkeiten
\newcommand{\MAN}{\textsc{Man}}
% \DeclareMathOperator{\GRUPPEN}{GRUPPEN} %Gruppen
\newcommand{\GRUPPEN}{\textsc{Gruppen}}
% \DeclareMathOperator{\ABELGRUPPEN}{ABEL.GRUPPEN} %abelsche Gruppen
\newcommand{\ABELGRUPPEN}{\textsc{Abel.Gruppen}}
% \DeclareMathOperator{\KAT}{KAT} %Kategorien
\newcommand{\KAT}{\textsc{Kat}}
% \DeclareMathOperator{\FUN}{FUN} %Funktoren
\newcommand{\FUN}{\textsc{Fun}}


\newcommand{\op}{\mathrm{op}}

\newcommand{\Ker}{\ensuremath{\Kern \,}}
\newcommand{\zirkel}{\rotatebox[origin = c]{-90}{$\varangle$}}
\newcommand{\wirkung}{\!\curvearrowright\!}

\newcommand{\oE}{\OE}

%--Beweisende
\newcommand{\bewende}{\ifmmode \tag*{$\square$} \else \hfill $\square$ \fi}

%--Volumen
\newcommand{\vol}[1]{\ensuremath{\Vol_{#1}}}

%--Skalarprodukt
\DeclarePairedDelimiterX\sprod[2]{\langle}{\rangle}{#1\,\delimsize\vert\,#2}

%--Betrag, Gaußklammer
\DeclarePairedDelimiter{\abs}{\lvert}{\rvert}
\DeclarePairedDelimiter{\ceil}{\lceil}{\rceil}

%--Norm
\DeclarePairedDelimiter\doppelstrich{\Vert}{\Vert}
\newcommand{\norm}[2][\relax]{
\ifx#1\relax \ensuremath{\doppelstrich*{#2}}
\else \ensuremath{\doppelstrich*{#2}_{#1}}
\fi}


%--Umklammern
\DeclarePairedDelimiter\enbrace{(}{)}
\DeclarePairedDelimiter\benbrace{[}{]}
\DeclarePairedDelimiter\lenbrace{<}{>}

%--Mengen
\DeclarePairedDelimiterX\mengenA[1]{\lbrace}{\rbrace}{#1}
\DeclarePairedDelimiterX\mengenB[2]{\lbrace}{\rbrace}{#1\, \delimsize\vert \, #2}

\makeatletter
\newcommand{\set}[2][\relax]{
\ifx#1\relax \ensuremath{
\mengenA*{#2}}
\else \ensuremath{%
  \mengenB*{#1}{#2}}
\fi}
\makeatother


%--offen und abgeschlossen
\newcommand{\off}{\! \stackrel[\text{offen}]{}{\subset} \!}
\newcommand{\abg}{\! \stackrel[\text{abg.}]{}{\subset} \!}

%--Differential
\newcommand{\diff}[2]{\ensuremath{\frac{{\partial #1}}{{\partial #2}} }}
\newcommand{\diffd}[2]{\ensuremath{\frac{\mathd #1}{\mathd #2} }}

%--Diagonale Linien für Matrizen
\newcommand\tikzmark[1]{%
  \tikzexternaldisable\tikz[overlay,remember picture,baseline] \node [anchor=base] (#1) {};\tikzexternalenable}

\newcommand\MyLine[3][]{%
  \tikzexternaldisable\begin{tikzpicture}[overlay,remember picture]
    \draw[#1] (#2.south) -- (#3.north);
  \end{tikzpicture}\tikzexternalenable}
  
%--Jordankasten
\newcommand{\Jord}{ \tikzexternaldisable \begin{BMAT}{|ccc|}{|ccc|} %
				0 & & \\ %
				1\tikzmark{z} &  \diagdown & \\ %
				& \tikzmark{u}1 \MyLine[thick]{z}{u} &  0 %
			\end{BMAT} \tikzexternalenable}
			
%--Jordankasten mit Lambda
\newcommand{\JordLa}[1][\empty]{%
	\tikzexternaldisable\ifthenelse{\equal{#1}{\empty}}
		{\begin{BMAT}{|ccc|}{|ccc|} %
				\lambda & & \\ %
				1\tikzmark{z} &  \diagdown & \\ %
				& \tikzmark{u}1 \MyLine[thick]{z}{u} &  \lambda  %
			\end{BMAT}}
		{\begin{BMAT}{|ccc|}{|ccc|} %
				\lambda_{#1}  & & \\ %
				1\tikzmark{z} &  \diagdown & \\ %
				& \tikzmark{u}1 \MyLine[thick]{z}{u} &  \lambda_{#1}  %
			\end{BMAT}}\tikzexternalenable}

%--Inhaltsverzeichnis
\usepackage[tocindentauto]{tocstyle}
\usetocstyle{KOMAlike}			

\newcommand{\fach}{Einführung in die Funktionalanalysis}
\newcommand{\semester}{WiSe 2014}
\newcommand{\homepage}{https://wwwmath.uni-muenster.de/u/wilhelm.winter/wwinter/funktionalanalysis.html}
\newcommand{\prof}{Prof.\,Dr.\,Wilhelm Winter}
\providecommand{\verfasser}{Jannes Bantje}

%--Konfiguration von scrheadings
\setheadsepline{1pt}[\color{light_gray}]
\pagestyle{scrheadings}
%\fancyhf[H, F]{}
\clearscrheadfoot

\providecommand{\shortFach}{\fach}
\lehead{ \includegraphics[height=0.5 cm,keepaspectratio]{../!config/Bilder/Logo_WWU_Muenster_light_gray.pdf}}
\rehead{\rule{0cm}{0.5cm}\footnotesize \sffamily \color{light_gray} \verfasser -- Mitschrift \shortFach}
\lohead{\rule{0cm}{0.5cm} \footnotesize \sffamily \color{light_gray} Stand: \today \; \thistime[:]}
\rohead{\includegraphics[height=0.5 cm,keepaspectratio]{../!config/Bilder/fb10logo_gray.pdf}}


\ofoot[{ \color{dark_gray} \LARGE \sffamily \thepage}]{{ \color{dark_gray} \LARGE \sffamily \thepage}} %hier wir auch der plain Stil bearbeitet!
\automark{section}
\ifoot{ \color{dark_gray} \small \leftmark}

%--Metadaten
\author{\verfasser}
\titlehead{\includegraphics[height=1.5cm, keepaspectratio]{../!config/Bilder/Logo_WWU_Muenster.pdf}%
\hfill \includegraphics[height=1.3cm, keepaspectratio]{../!config/Bilder/fb10logo.pdf}}
\title{Skript \fach}
\subtitle{Mitschrift der Vorlesung  \enquote{\fach} von \prof}

\publishers{
Aktuelle Version verfügbar bei: \bigskip\\
\normalsize
\begin{minipage}[t]{0.45\textwidth}
	\centering{
	\qrcode[height=3.5cm, version=4]{https://github.com/JaMeZ-B/latex-wwu}
	\medskip\\
	\includegraphics[height=0.4cm, keepaspectratio]{../!config/Bilder/github_octo.pdf} 
	\includegraphics[height=0.4cm, keepaspectratio]{../!config/Bilder/GitHub_Logo.pdf} (inklusive Sourcecode)\\
	{\small\url{https://github.com/JaMeZ-B/latex-wwu}}}
\end{minipage}
\hfill
\begin{minipage}[t]{0.45\textwidth}
	\centering{
	\qrcode[height=3.5cm, version=4]{btsync://B6WH2DISQ5QVYIRYIEZSF4ZR2IDVKPN3I?n=latex_share}
	\medskip\\
	\includegraphics[height=0.4cm, keepaspectratio]{../!config/Bilder/bt_sync_logo.pdf} 
	{\large \textbf{Bittorrent} Sync} \\
	\texttt{B6WH2DISQ5QVYIRYIEZSF4ZR2IDVKPN3I}
	}
\end{minipage}
}



\begin{document}
\pagenumbering{Roman}
\maketitle
\begin{abstract}
\section*{Vorwort --- Mitarbeit am Skript}
Dieses Dokument ist eine Mitschrift aus der Vorlesung \enquote{\fach, \semester}, gelesen von \prof. Der Inhalt entspricht weitestgehend dem Tafelanschrieb. Für die
Korrektheit des Inhalts übernehme ich keinerlei Garantie! Für Bemerkungen und Korrekturen -- und seien es nur Rechtschreibfehler -- bin ich sehr dankbar. 
Korrekturen lassen sich prinzipiell auf drei Wegen einreichen: 
\begin{itemize}
	\item Persönliches Ansprechen in der Uni, Mails an \hrefsymmail{mailto:\mail}{\mail} (gerne auch mit annotieren PDFs) 
	\item \emph{Direktes} Mitarbeiten am Skript: Den Quellcode poste ich auf GitHub (siehe oben), also stehen vielfältige Möglichkeiten der Zusammenarbeit zur Verfügung:
	Zum Beispiel durch Kommentare am Code über die Website und die Kombination Fork + Pull Request. Wer sich verdient macht oder ein Skript zu einer Vorlesung, die 
	ich nicht besuche, beisteuern will, dem gewähre ich gerne auch Schreibzugriff.
	
	Beachten sollte man dabei, dass dazu ein Account bei \url{github.com} notwendig ist, der allerdings ohne Angabe von persönlichen Daten angelegt werden kann. 
	Wer bei GitHub (bzw. dem zugrunde liegenden Open-Source-Programm \enquote{\texttt{git}}) -- verständlicherweise -- Hilfe beim Einstieg braucht, dem helfe ich gerne 
	weiter. Es gibt aber auch zahlreiche empfehlenswerte Tutorials im Internet.\footnote{zB. \url{https://try.github.io/levels/1/challenges/1}, ist auf Englisch, aber dafür 
	interaktives LearningByDoing}
	\item \emph{Indirektes} Mitarbeiten: \TeX-Dateien per Mail verschicken. 
	
	Dies ist nur dann sinnvoll, wenn man einen ganzen Abschnitt ändern möchte (zB. einen alternativen Beweis geben), da ich die Änderungen dann per Hand einbauen muss! Ich freue mich aber auch über solche Beiträge!
\end{itemize}
\section*{Vorlesungshomepage}
{\centering 
\begin{minipage}[c][][c]{\textwidth}
	\centering \qrcode[height=3cm]{\homepage} \medskip\\
	\footnotesize \url{\homepage}
\end{minipage}
\par}

\end{abstract}

\setcounter{page}{1}
\tableofcontents
\cleardoubleoddemptypage
\pagenumbering{arabic}
\setcounter{page}{1}

\section{Metrische Räume und der Satz von Baire} % (fold)
\label{sec:1}

\subsection[Definition: Metrischer Raum]{Definition} % (fold)
\label{sub:11}
Ein \Index{metrischer Raum} ist ein Paar $(X,d)$, wobei $X$ eine Menge und $d : X \times X \to [0,\infty)$ ist, sodass\marginnote{14 Okt}
\begin{enumerate}[1)]
	\item $d(x,y)=0 \iff x=y$ $\forall x,y \in X$
	\item $d(x,y) = d(y,x)$ $\forall x,y \in X$
	\item $d(x,z) \le d(x,y) + d(y,z)$ $\forall x,y,z \in X$
\end{enumerate}
% subsection 11 (end)

\subsection[Definition: Offen, abgeschlossen und Abschluss]{Definition} % (fold)
\label{sub:12}
Sei $(X,d)$ ein metrischer Raum
\begin{itemize}
	\item Eine Teilmenge $U \subseteq X$ heißt \Index{offen}, falls für jedes $x \in U$ ein $\varepsilon > 0$ existiert, so dass 
	\[
		B(x,\varepsilon) := \set[y \in X]{d(x,y) < \varepsilon} \subset U 
	\]
	\item Eine Teilmenge $A \subset X$ heißt \Index{abgeschlossen}, falls $X \setminus A$ offen ist (als Teilmenge von $X$). 
	\item $\mathcal{T}X := \set[U \subset X]{U \text{ offen}} $ ist die Topologie auf $X$ (die von der Metrik $d$ induziert wird)
	\item Falls $W \subset X$ eine Teilmenge ist, dann bezeichnet $\overline{W}$ den \Index{Abschluss} von $W$, d.h. die kleinste abgeschlossene Teilmenge von $X$, die 
	$W$ enthält. Es gilt
	\[
		\overline{W} = \enspace\bigcap_{\mathclap{A \subset X \text{abg.}, W \subset A}} \enspace A 
	\] 
\end{itemize}
Für metrische Räume gilt:
\(
	\overline{W} = \set[\lim_n x_n]{(x_n)_n \subset W \text{ konvergente Folge}}.
\)
Warum gilt $W \subset \overline{W}$?
% subsection 12 (end)

\subsection[Definition: Stetigkeit, gleichmäßige Stetigkeit, Isometrie]{Definition} % (fold)
\label{sub:13}
Eine Abbildung $f : X \to Y$ zwischen zwei metrischen Räumen $(X,d_X)$ und $(Y,d_Y)$ heißt 
\begin{itemize}
	\item \bet{stetig in} $x \in X$, falls
	\(
		\forall \varepsilon >0 : \exists \delta >0 : \forall x' \in X : d_X(x,x') < \delta  \Longrightarrow d_Y \enbrace[\big]{f(x), f(x')} < \varepsilon 
	\)
	\item \Index{stetig}, falls $f$ an jedem Punkt $x \in X$ stetig ist.
	
	(Äquivalent: Für jede offene Menge $V$ in $Y$ ist $f ^{-1}(V)$ offen in $X$) 
	\item \Index{gleichmäßig stetig}, falls
	\(
		\forall \varepsilon>0 : \exists \delta >0 : \forall x,x' \in X : d_X(x,x') < \delta \Longrightarrow d_Y \enbrace[\big]{f(x), f(x')} < \varepsilon .
	\) 
	\item $f$ heißt \Index{Isometrie}, falls \marginnote{automatisch injektiv}
	\(
		\forall x,x' \in X : d_X(x,x') = d_Y \enbrace[\big]{f(x), f(x')}.
	\)
	\item $f$ heißt \Index{isometrischer Isomorphismus}, falls $f$ bijektiv und isometrisch ist. \marginnote{$f ^{-1}$ auch}
\end{itemize}
% subsection 13 (end)

\subsection[Definition: Cauchy-Folge und Vollständigkeit]{Definition} % (fold)
\label{sub:14}
Eine Folge $(x_n)_{n \in \mathds{N}}$ in einem metrischen Raum $(X,d)$ heißt \bet{Cauchy}\index{Cauchy-Folge}, falls
\[
	\forall \varepsilon>0 : \exists N \in \mathds{N} : \forall n,k >N : d(x_n,x_k) < \varepsilon.
\]
$(X,d)$ heißt \Index{vollständig}, falls jede Cauchy-Folge in $X$ konvergiert. 
% subsection 14 (end)

\subsection[Satz: Existenz einer eindeutigen Vervollständigung metrischer Räume]{Satz} % (fold)
\label{sub:15}
Sei $(X,d)$ ein metrischer Raum. Dann existiert ein vollständiger metrischer Raum $(\tilde{X}, \tilde{d})$ und eine Isometrie $\iota : X \hookrightarrow \tilde{X}$, sodass
$\overline{\iota(X)} = \tilde{X}$ (d.h. $\iota(X)$ ist dicht in $\tilde{X}$). $(\tilde{X}, \tilde{d})$ heißt \Index{Vervollständigung} von $(X,d)$ und ist eindeutig bis auf
isometrische Isomorphie.
\minisec{Beweis}
\begin{description}
	\item[Eindeutigkeit:] 
	Angenommen, $(\hat{X}, \hat{d})$ ist ein weiterer vollständiger metrischer Raum und $\kappa : X \to \hat{X}$ eine Isometrie mit 
	$\overline{\kappa(X)}=\hat{X}$. Definiere $\gamma : \hat{X} \to \tilde{X}$ wie folgt: Falls $y \in \hat{X}$, wähle eine Folge $(x_n)_{n \in \mathds{N}}$ in $X$, sodass
	$y= \lim_{n \to \infty} \kappa(x_n)$. Setze nun
	\[
		\gamma(y) := \lim_{ n \to \infty} \iota(x_n) \in \tilde{X} 
	\]
	Zu zeigen: $\gamma$ ist ein isometrischer Isomorphismus.\todo{Achtung Wortspiel: Beweis \emph{vervollständigen} ;-D}
	\begin{description}
		\item[Injektivität:] Seien $y,y' \in \hat{X}$ mit $\gamma(y)=\gamma(y')$. Dann existieren Folgen $(x_n)_n, (x_n')_n$ in $X$ mit
		\[
			\lim_{ n \to \infty} \iota(x_n) = \lim_{ n \to \infty} \iota(x_n') 
		\]
	\end{description}
	\item[Existenz:] \begin{description}
		\item[Konstruktion von $(\tilde{X}, \tilde{d})$:] Setze $Y := \set[(x_n)_{n \in \mathds{N}}]{(x_n)_n \text{ ist Cauchy-Folge in }X}$. Definiere 
	\[
		(x_n)_n \sim (x_n') :\iff \lim_{ n \to \infty} d(x_n, x_n') = 0
	\]
	$\sim$ ist eine Äquivalenzrelation auf $Y$. Definiere nun $\tilde{X} := \nicefrac{Y}{\sim}$ und $\tilde{d} :  \tilde{X} \times \tilde{X} \to [0, \infty)$ durch
	\[
		\tilde{d} \enbrace*{\benbrace[\big]{(x_n)_n} , \benbrace[\big]{(x_n')_n}  } := \lim_{ n \to \infty} d(x_n , x_n') 
	\]
	$\tilde{d}$ ist eine wohldefiniert, d.h. falls $\benbrace*{(x_n)_n} = \benbrace*{(y_n)_n} $ und $\benbrace*{(x_n')_n} = \benbrace*{(y_n')_n}$, dann ist
	\[
		\lim_{n \to \infty} d(x_n , x_n') = \lim_{n \to \infty} d(y_n, y_n'). \tag*{(leichte Übung)}
	\]
	Weiter ist $\enbrace[\big]{d(x_n x_n')}_n$ eine Cauchy-Folge in dem vollständigen Raum $[0,\infty)$ und somit 
	konvergent: Sei dazu $\varepsilon>0$. Dann existiert ein $N \in \mathds{N}$, sodass $d(x_n,x_m) \le \frac{\varepsilon}{2}$ und $d(x_n', x_m') \le \frac{\varepsilon}{2}$ für alle $n,m \ge N$. Dann gilt nach der \hyperref[sub:vier_ungl]{Vierecksungleichung (siehe Anhang \ref*{sub:vier_ungl})}
	\[
		\abs[\big]{d(x_n,x_n') - d(x_m,x_m')} \le d(x_n, x_m) + d(x_n', x_m') \le \varepsilon
	\]
	\item[Einbettung von $X:$] Definiere nun $\iota : X \to \tilde{X}$ durch $x \mapsto [(x,x,x, \ldots )] \in \tilde{X}$. $\iota$ ist Isometrie, da
	\[
		\tilde{d} \enbrace[\big]{\iota(x),\iota(y)} = \lim_{ n \to \infty} d(x,y) = d(x,y)
	\]
	% Falls $\benbrace*{(x_n)_n} \in \tilde{X}$, finde $y_k \in \iota(X)$, so dass $\benbrace*{(x_n)} = \lim_{ k} y_k$ und $y_k = \benbrace*{(x_k, x_k, \ldots )}= \iota(x_k)$.
% 	Dann gilt $y_k \xrightarrow{k \to \infty} \benbrace*{(x_n)_n} $, d.h.
% 	\[
% 		\lim_{ n \to \infty} d(x_k,x_n) = \tilde{d} \enbrace*{y_k, \benbrace*{(x_n)_n}} \xrightarrow{k \to \infty} 0
% 	\]
% 	$\leadsto \tilde{X}= \overline{\iota(X)}$.
	Sei nun $\benbrace*{(x_n)_n} \in \tilde{X}$ und $\varepsilon>0$. Da $(x_n)_n$ eine Cauchy-Folge ist, gibt es ein $N \in \mathds{N}$, sodass für alle
	$n,m \ge N$ gilt $d(x_n,x_m) <\varepsilon$. Dann gilt
	\[
		\tilde{d}\enbrace[\big]{\iota(x_N), \benbrace*{(x_n)_n} } = \lim_{ n \to \infty} d \enbrace*{x_N, x_n} < \varepsilon
	\]
	\item[Vollständigkeit von $(\tilde{X}, \tilde{d})$:] Sei $\enbrace*{\overline{x}^m}_m$ eine Cauchyfolge in $\tilde{X}$.\marginnote{$\overline{x} \in \tilde{X}$}
	\[
		\Longrightarrow \forall \varepsilon >0 : \exists M(\varepsilon) \in \mathds{N} : \forall m,m' > M : \tilde{d} \enbrace*{\overline{x}^m, \overline{x}^{m'}} < \frac{\varepsilon}{3} 
	\]
	Wenn $\overline{x}^m = \benbrace*{(x^m_n)_n}$ und $\overline{x}^{m'} = \benbrace*{(x^{m'}_n)_n} $, dann gilt also für alle $m,m' >M(\varepsilon)$
	\begin{equation*}
		\lim_{ n \to \infty} d \enbrace*{x^m_n, x^{m'}_n} < \frac{\varepsilon}{3}  \label{eq:15:absch1}\tag{\dag}
	\end{equation*}
	Für alle $m$ ist $(x^m_n)_n$ eine Cauchyfolge, also gilt
	\begin{equation*}
		\forall m : \exists N(m) : \forall n,n' \ge N(m) : d \enbrace*{x^m_n, x^m_{n'}} < \frac{1}{m}  \label{eq:15:absch2}\tag{\dag \dag}
	\end{equation*}
	Setze nun $z_n := x^n_{N(n)}$. Behauptung: $(z_n)_n$ ist eine Cauchyfolge. Sei $\varepsilon>0$. Dann gilt für $n,m>N=\max\set{M(\varepsilon), \frac{3}{\varepsilon}}$
	für ein $k > N(m), N(n)$ \marginnote{$d(z_n,z_m)$ ist unabhängig von $k$, also kann man $k$ beliebig groß wählen}
	\begin{align*}
		d(z_n,z_m) = d \enbrace*{x^n_{N(n)}, x^m_{N(m)}} \le 
		\underbrace{d\enbrace*{x^n_{N(n)}, x^n_k}}_{\stackrel{\text{\eqref{eq:15:absch2}}}{<} \frac{1}{n} < \frac{1}{N} \le \frac{\varepsilon}{3}} + 
		\underbrace{d\enbrace[\Big]{x^n_k, x^m_k}}_{\stackrel{\text{\eqref{eq:15:absch1}}}{\le} \frac{\varepsilon}{3}}   + 
		\underbrace{d\enbrace*{x^m_k, x^m_{N(m)}}}_{\stackrel{\text{\eqref{eq:15:absch2}}}{<} \frac{1}{m} < \frac{1}{N} \le \frac{\varepsilon}{3}} 
		< \varepsilon
	\end{align*}
	$\Rightarrow (z_n)_n$ ist eine Cauchyfolge, also 
	\begin{equation*}
		\forall \varepsilon>0 : \exists N_z(\varepsilon) : \forall n,m > N_z(\varepsilon)  : d(z_n,z_m) < \varepsilon  \label{eq:15:absch3} \tag{\#}
	\end{equation*}
	Es bleibt zu zeigen: $\lim_{m \to \infty} \overline{x}^m = \benbrace*{(z_n)_n}$. Sei dazu $\varepsilon>0$. Dann gilt für 
	$m \ge \max \set{\frac{2}{\varepsilon}, N_z(\varepsilon)}$
	\begin{align*}
		\tilde{d} \enbrace[\big]{\overline{x}^m, \benbrace*{(z_n)_n}} = \lim_{ n \to \infty} d \enbrace*{x^m_n, x^n_{N(n)}}\le \lim_{ n \to \infty} 
		\enbrace[\Bigg]{ \underbrace{d\enbrace*{x_n^m, x^m_{N(m)}}}_{\stackrel{\text{\eqref{eq:15:absch2}}}{<} \frac{1}{m} \le \frac{\varepsilon}{2}  } +
		\underbrace{d\enbrace*{x^m_{N(m)}, x^n_{N(n)}}}_{=d(z_m,z_n) \stackrel{\text{\eqref{eq:15:absch3}}}{<} \frac{\varepsilon}{2} } } < \varepsilon
	\end{align*}
	Also gilt $\overline{x}^m \xrightarrow{m \to \infty} \benbrace*{(z_n)_n}$ und $(\tilde{X},\tilde{d})$ ist vollständig. \bewende
	\end{description}
\end{description}
% subsection 15 (end)

\subsection[Definition: Raum der beschränkten, stetigen Abbildungen]{Definition} % (fold)
\label{sub:16}
Sei $(W,\mathcal{T})$ ein topologischer Raum und $(X,d)$ ein metrischer Raum. Sei 
\[
	C_b(W,X) = \set[f : W \to X]{f \text{ stetig und beschränkt}} 
\]
versehen mit der Metrik $d_{W,X}$, definiert durch
\[
	d_{W,X} = \sup_{t \in W} d \enbrace[\big]{f(t), g(t)} 
\]
% subsection 16 (end)

\subsection[Bemerkung: $d_{W,X}$ als Metrik auf $C(W,X)$]{Bemerkung} % (fold)
\label{sub:17}
Auf $C(W,X) = \set{f : W \to X \text{ stetig}} $ ist $d_{W,X}$ eine \Index{erweiterte Metrik}, d.h. der Wert $\infty$ ist möglich.
$\tilde{d}_{W,X} := \min \set{1, d_{W,X}} $ ist eine \enquote{echte} Metrik auf $C(W,X)$.
% subsection 17 (end)

\subsection[Proposition: $X$ vollständig $\Rightarrow$ $C-B(W,X)$ und $C8W,X$ vollständig]{Proposition} % (fold)
\label{sub:18}
Falls $X$ vollständig ist, dann sind $C_b(W,X)$ und $C(W,X)$ vollständig (bezüglich $d_{W,X}$ bzw. $\tilde{d}_{W,X}$).
\minisec{Beweis}
Übung.\todo{nach Übung hinzufügen}
% subsection 18 (end)

\subsection[Proposition über eine Folge von Bällen]{Proposition} % (fold)
\label{sub:19}
Sei $(X,d)$ ein vollständiger metrischer Raum und $\enbrace*{\overline{B}(x_n, \varepsilon_n)}_{n \in \mathds{N}}$ mit 
$\overline{B}(x_{n+1}, \varepsilon_{n+1}) \subseteq \overline{B}(x_n,\varepsilon_n)$ und $\varepsilon_n \xrightarrow{n \to \infty} 0$. Dann existiert genau 
ein Punkt in $\bigcap_{n \in \mathds{N}} \overline{B}(x_n, \varepsilon_n)$.
\minisec{Beweis}
\begin{description}
	\item[Existenz:] Behauptung: Die Folge der Mittelpunkte $(x_n)_{n \in\mathds{N}}$ ist eine Cauchy-Folge: Sei $\varepsilon >0$. Finde $N \in \mathds{N}$,
	sodass $\varepsilon_n \le \varepsilon$ für alle $n \ge N$. Dann gilt für alle $n \ge m \ge N$
	\[
		d(x_n, x_m) \le \varepsilon_m \le \varepsilon,
	\]
	da $x_n \in \overline{B}(x_m, \varepsilon)$ ist. Da $X$ vollständig ist, existiert ein $x \in X$ mit $x_n \xrightarrow{n \to \infty} x$.
	
	Behauptung: $x \in \bigcap_{n \in \mathds{N}} \overline{B}(x_n, \varepsilon_n)$. Wähle dazu ein $N \in \mathds{N}$. Dann ist $(x_n)_{n \ge N}$ eine
	Cauchy-Folge in $\overline{B}(x_N,\varepsilon_N)$. $\overline{B}(x_N,\varepsilon_N)$ ist ein abgeschlossener Teilraum von $X$ und somit vollständig. 
	Also ist $x \in \overline{B}(x_N,\varepsilon_N)$. Da $N$ beliebig war, gilt $x \in \overline{B}(x_n,\varepsilon_n)$ für jedes $n \in \mathds{N}$.
	\item[Eindeutigkeit:] Es seien $x,y \in \bigcap_{n \in \mathds{N}} \overline{B}(x_n,\varepsilon_n)$, also 
	$x,y \in \overline{B}(x_n, \varepsilon_n) $ für alle $ n \in \mathds{N}$. Dann gilt
	\[
		d(x,y) \le d(x,x_n) + d(x_n,y) \le 2 \cdot \varepsilon_n \xrightarrow{n \to \infty} 0
	\]
	für jedes $n \in \mathds{N}$ und somit $d(x,y)=0 \Leftrightarrow x=y$. \bewende
\end{description}
% subsection 19 (end)
% section 1 (end)




\cleardoubleoddemptypage
\appendix
\section{Anhang} % (fold)
\label{sec:anhang}

\subsection{Vierecksungleichung} % (fold)
\label{sub:vier_ungl}
Sei $(X,d)$ ein metrischer Raum und $x,y,u,v \in X$. Dann gilt
\[
	\abs[\big]{d(x,y) - d(u,v)} \le d(x,u) + d(y,v) 
\]
\minisec{Beweis}
Einerseits gilt nach der Dreiecksungleichung
\[
	d(x,y) \le d(x,u) + d(u,v) + d(v,y) \quad \Longrightarrow \quad d(x,y) - d(u,v) \le d(x,u) + d(y,v)
\]
Andererseits aber auch
\[
	d(u,v) \le d(u,x) + d(x,y) + d(y,v) \quad \Longrightarrow \quad d(u,v) - d(x,y) \le d(x,u) + d(y,v)
\]
Insgesamt folgt also die Behauptung. \bewende
% subsection vier_ungl (end)

% section anhang (end)
\cleardoubleoddemptypage
\pagenumbering{Alph}
\setcounter{page}{1}
\printindex
\listoffigures
\end{document}
